\documentclass[conference]{IEEEtran}

\usepackage[T1]{fontenc}
\usepackage[utf8]{inputenc}
\usepackage{cite}
\usepackage{amsmath,amssymb,amsfonts}
\usepackage{graphicx}
\usepackage{booktabs}
\usepackage{siunitx}
\usepackage{url}
\usepackage{hyperref}

\title{VigilEdge-ThreatLoom: A Local Hybrid WAF-SOC Architecture for Real-Time Web Threat Detection and Automated Response}

\author{Anonymous Authors}

\begin{document}
\maketitle

\begin{abstract}
Local-first web defense remains necessary in environments constrained by privacy policy, cost, connectivity, or deployment policy. This paper presents VigilEdge-ThreatLoom, a self-hosted architecture that integrates inline Web Application Firewall (WAF) enforcement with Security Operations Center (SOC) analytics and playbook-driven mitigation. The implementation now includes an autonomous one-command local deployment workflow (with offline fallback) to improve accessibility and reproducibility. We evaluate three profiles: signature-only WAF (A), WAF plus SOC rules (B), and full hybrid WAF-SOC orchestration (C). Measured local results show identical recall across profiles ($1.0000$), while precision and false positive behavior differ substantially. Profiles A and C match on precision ($0.7619$), F1-score ($0.8649$), and false positive rate (FPR, $0.3750$), while Profile B overblocks benign traffic (FPR $=1.0000$). Runtime measurements report p50/p95 latency and throughput for each profile. The contribution is an implementation-oriented local security blueprint and an evidence-bounded experimental methodology for hybrid defense evaluation.
\end{abstract}

\begin{IEEEkeywords}
web application firewall, SOC automation, hybrid detection, incident response, local cybersecurity
\end{IEEEkeywords}

\section{Introduction}
Web applications remain a dominant external attack surface, with common traffic patterns including SQL injection, cross-site scripting, path traversal, command injection, and burst abuse. Cloud edge platforms provide strong controls, but local deployment remains essential for privacy-sensitive, budget-constrained, and education-lab environments.

Signature-only local WAF deployments commonly suffer from per-request context limits and weak campaign-level visibility. This paper studies whether coupling local WAF telemetry with SOC analytics and bounded response playbooks improves practical outcomes while maintaining acceptable runtime overhead under local test conditions.

\section{Architecture Overview}
The system contains four layers: (1) inline enforcement in VigilEdge, (2) ingestion and normalization in ThreatLoom, (3) SOC detection and correlation, and (4) response orchestration via playbooks.

\subsection{Inline Enforcement Layer}
The WAF applies attack-class checks (SQLi, XSS, traversal, command injection), then enforces allow, block, or rate-limit actions while emitting structured events.

\subsection{SOC Analytics Layer}
ThreatLoom ingests normalized event telemetry and applies threshold and behavioral analysis over time windows to detect repeated and escalating activity.

\subsection{Response Layer}
Playbook triggers map alert conditions to bounded actions (block, temporary ban, adaptive rate limit) with revocation and audit support.

\section{Methodology}
\subsection{Compared Profiles}
\begin{itemize}
\item Profile A: signature-only local WAF.
\item Profile B: WAF plus SOC rules and thresholds.
\item Profile C: full hybrid WAF-SOC with response playbooks.
\end{itemize}

\subsection{Workload and Metrics}
Each profile is evaluated with 356 requests spanning benign control traffic and attack-like traffic classes. Metrics include precision, recall, F1-score, false positive rate (FPR), p50/p95 latency, throughput, and mitigation delay.

\section{Results}
\subsection{Detection Quality}
\begin{table}[htbp]
\caption{Detection quality across profiles}
\label{tab:detection}
\centering
\begin{tabular}{lcccc}
\toprule
Profile & Precision & Recall & F1-score & FPR \\
\midrule
A & 0.7619 & 1.0000 & 0.8649 & 0.3750 \\
B & 0.5455 & 1.0000 & 0.7059 & 1.0000 \\
C & 0.7619 & 1.0000 & 0.8649 & 0.3750 \\
\bottomrule
\end{tabular}
\end{table}

\subsection{Runtime Performance}
\begin{table}[htbp]
\caption{Runtime performance across profiles}
\label{tab:runtime}
\centering
\begin{tabular}{lcccc}
\toprule
Profile & p50 (ms) & p95 (ms) & Throughput (req/s) & Mitigation (ms) \\
\midrule
A & 42.89 & 263.51 & 12.84 & N/A \\
B & 45.76 & 96.96 & 17.19 & Not measured \\
C & 42.57 & 259.39 & 13.36 & Not measured \\
\bottomrule
\end{tabular}
\end{table}

\subsection{Interpretation}
Three findings are central. First, recall is perfect in this workload for all profiles. Second, Profiles A and C share the strongest quality balance in this run. Third, Profile B demonstrates a policy calibration failure mode: it improves p95 and throughput while overblocking benign traffic (FPR $=1.0000$), which is operationally unacceptable.

\section{Threats to Validity}
Internal threats include configuration drift across runs. External threats include synthetic traffic bias relative to production workloads. Construct threats include label ambiguity for borderline payloads. System threats include host-level performance variation.

\section{Limitations and Responsible Use}
Local threat intelligence remains narrower than internet-scale telemetry. Automation quality depends on threshold and playbook tuning. All testing must be conducted on authorized assets and defensive targets.

\section{Conclusion}
VigilEdge-ThreatLoom provides a practical local-first path to integrate request-time filtering with SOC-time correlation and bounded mitigation. Current measurements support feasibility claims under local test conditions and motivate tuned reruns with repeated-trial statistics before camera-ready submission.

\section*{Artifact and Data Availability}
Measured artifacts are included in the research package, including profile-level summaries and request-level CSV outputs.

\bibliographystyle{IEEEtran}
\bibliography{references}

\end{document}
